\scanpage{52}
Pour ceci, il faut supposer en plus des
conditions a) b) de 1'), les conditions
\[
(21)\quad\left\{\begin{aligned}
\text{c)}\quad&X=\bigcup X_i\\
\text{d)}\quad&\forall(i,j)\in\unsure{I|X}\times I,\ \text{on a}\quad
X_i\cap X_j=\bigcup_{\substack{k\in I\\k\leq i,\ k\leq j}}X_k,
\end{aligned}\right.
\]
qui assurent trivialement la relation (20).
\deleted{On devrait \unsure{pouvoir en déduire}}\inserted{C'est équivalent}, de façon \deleted{\unsure{e}} ``formelle'' :

\uline{Proposition} Sous les conditions a) b) c) d)
de (1') (21), l'application canonique
\[
(22)\qquad\underset{\longrightarrow}{\mathrm{Lim}}\ \underbrace{X_{\Delta_*}}\xrightarrow{\sim}X
\]
\diagram{\begin{tikzcd}[column sep=small]
\{X_{\Delta_0}&X_{\Delta_1}\arrow[l,shift left=1ex]\arrow[l,shift right=1ex]&X_{\Delta_2}\arrow[l,shift left=1.5ex]\arrow[l]\arrow[l,shift right=1.5ex]\ \ldots\}
\end{tikzcd}}
est un homéomorphisme.

Cette proposition est valable avec les
deux variantes (\deleted{simplexes}\inserted{drapeaux} quelconques, dra-
peaux non dégénérés) pour les\unsure{.}
espaces de drapeaux $X_{\Delta_r}$.

\uline{Corollaire} Soit : $(X_{\Delta_0},X_{\Delta_1}\subset X_{\Delta_0}\times X_{\Delta_0})$
un espace topologique ordonné, \deleted{supposons
$X_{\Delta_1}$ \unsure{est une partie fermée} de $X_{\Delta_0}\times X_{\Delta_0}$.}
\inserted{et soit $(X_{\deleted{\unsure{\Delta_0|\Delta}}i})_{i\in I}$ une \deleted{\unsure{famille}} partition de}
\deleted{\unsure{supposons}} \deleted{\unsure{où}} $X_{\Delta_0}$ en des parties à la fois ouvertes et fermées
\margin{51}
